% !TEX program = pdflatex
% Supporting Information for Paper 2 (Molecular Informatics).
% Compile twice for cross-refs.
\documentclass[11pt,a4paper]{article}
\usepackage[T1]{fontenc}\usepackage[utf8]{inputenc}\usepackage{lmodern}
\usepackage[margin=1in]{geometry}
\usepackage{amsmath,amssymb,amsthm}\usepackage{mathtools}
\usepackage{booktabs}\usepackage{array}\usepackage{graphicx}
\graphicspath{{../figures/}}
\usepackage{hyperref}\hypersetup{colorlinks=true,linkcolor=blue!50!black,citecolor=blue!50!black,urlcolor=blue!50!black}
\usepackage{xcolor}\usepackage{enumitem}
\newtheorem{theorem}{Theorem}\theoremstyle{definition}\newtheorem{remark}[theorem]{Remark}
\newcommand{\bid}{\mathcal{F}_{\mathrm{BID}}}
\renewcommand{\thesection}{S\arabic{section}}
\renewcommand{\thetable}{S\arabic{table}}
\title{\bfseries Supporting Information for:\\ Orthogonality Screening of Topological Indices for QSPR Modeling}
\author{G.~Shiwakoti, A.~Natarajan, P.~Chalise, M.~Arockiaraj}
\date{}
\begin{document}\maketitle

\section{Proof of the Edge-Degree-Pair Basis (Theorem~1 of the main text)}

\begin{theorem}[Edge-Degree-Pair Basis]\label{thm:bid}
Let $\mathcal{G}_\Delta$ denote the class of finite simple graphs
with maximum degree at most $\Delta$. For any BID index
$I_f\in\bid$ and any $G\in\mathcal{G}_\Delta$,
\begin{equation}\label{eq:bid_linear}
I_f(G)=\sum_{1\le i\le j\le\Delta} f(i,j)\cdot m_{ij}(G),
\end{equation}
where $m_{ij}(G)=\bigl|\{uv\in E(G):\{d_u,d_v\}=\{i,j\}\}\bigr|$.
Restricted to $\mathcal{G}_4$, the vector space spanned by all BID
indices has dimension at most $10$.
\end{theorem}

\begin{proof}
We establish the linear-decomposition identity~\eqref{eq:bid_linear}
first and then deduce the dimension bound.

\emph{Step 1: linear decomposition.} Fix $G\in\mathcal{G}_\Delta$ and,
for each pair $(i,j)$ with $1\le i\le j\le\Delta$, define
$E_{ij}(G) = \{ uv\in E(G) \,:\, \{d_u,d_v\}=\{i,j\}\}$,
so that $m_{ij}(G)=|E_{ij}(G)|$. Because every edge has
$d_u,d_v\in\{1,\ldots,\Delta\}$, the family $\{E_{ij}(G)\}_{i\le j}$
partitions $E(G)$. Hence
\[
I_f(G) = \sum_{uv\in E(G)} f(d_u,d_v)
= \sum_{1\le i\le j\le\Delta}\;\sum_{uv\in E_{ij}(G)} f(d_u,d_v).
\]
For every $uv\in E_{ij}(G)$ the multiset $\{d_u,d_v\}$ equals
$\{i,j\}$ and $f$ is symmetric, so $f(d_u,d_v)=f(i,j)$ is constant
on the inner sum, which collapses to $f(i,j)\cdot m_{ij}(G)$ and
yields~\eqref{eq:bid_linear}.

\emph{Step 2: dimension bound.} For each admissible pair $(i,j)$
with $1\le i\le j\le\Delta$, regard the edge-count
$m_{ij}:\mathcal{G}_\Delta\to\mathbb{N}_{\ge 0}$ as a functional on
$\mathcal{G}_\Delta$. Equation~\eqref{eq:bid_linear} shows
$I_f = \sum_{i\le j} f(i,j)\cdot m_{ij}$ as functionals on
$\mathcal{G}_\Delta$, so
$\bid \subseteq \operatorname{span}_{\mathbb{R}}\{m_{ij}:1\le i\le j\le\Delta\}$.
A subspace of a span of $N_\Delta$ vectors has dimension at most
$N_\Delta$, hence
\[
\dim_{\mathbb{R}}\bid \le N_\Delta = \frac{\Delta(\Delta+1)}{2}
\]
by the standard multiset-counting identity. For $\Delta=4$ this is
$N_4=10$.
\end{proof}

\section{The 30-index classical baseline}

\begin{table}[h]
\centering
\caption{$30$-index classical baseline used for orthogonality screening.}
\label{tab:baseline}
\small
\begin{tabular}{ll}
\toprule
Family & Indices\\
\midrule
Degree-based (18) & $M_1$, $M_2$, $mM_1$, $mM_2$, $F$, $R$, SCI, $H$, $GA$, $AG$, ABC, ABS, AZI, $SO$, $SO_{\mathrm{red}}$, Alb, $\sigma$, $\mathrm{redM}_1$\\
Distance-based (8) & $W$, WW, HR, MTI (Schultz), Gutman, Mostar, $Sz$, $PI$\\
Spectral (3) & EE, Energy, Spectral Radius\\
Other (1) & Balaban $J$\\
\bottomrule
\end{tabular}
\end{table}

\section{Threshold-sensitivity of the Wuzi screening verdict}

\begin{table}[!htbp]
\centering
\small
\caption{Wuzi $100$-point grid: pass counts under threshold variation. Top block: pairwise screen only ($\max_b\lvert r(W, X_b)\rvert < \tau_r$). Bottom block: combined diagnostic (pairwise pass \emph{and} $\lvert\mathrm{pcor}(W, y \mid X)\rvert \ge \tau_{\mathrm{pcor}}$) across the full $3 \times 3$ threshold grid.}
\label{tab:threshold_sensitivity}
\begin{tabular}{lccc}
\toprule
& \multicolumn{3}{c}{Pairwise pass count ($\tau_r$)}\\
\cmidrule(lr){2-4}
Dataset & $\tau_r = 0.90$ & $\tau_r = 0.95$ & $\tau_r = 0.99$\\
\midrule
ESOL          & $0$ & $0$ & $27$\\
FreeSolv      & $0$ & $1$ & $37$\\
Lipophilicity & $0$ & $0$ & $30$\\
\bottomrule
\end{tabular}

\vspace{0.6em}

\begin{tabular}{l|ccc|ccc|ccc}
\toprule
& \multicolumn{3}{c|}{$\tau_r = 0.90$} & \multicolumn{3}{c|}{$\tau_r = 0.95$} & \multicolumn{3}{c}{$\tau_r = 0.99$}\\
$\tau_{\mathrm{pcor}}$ & $.05$ & $.10$ & $.20$ & $.05$ & $.10$ & $.20$ & $.05$ & $.10$ & $.20$\\
\midrule
ESOL          & $0$ & $0$ & $0$ & $0$ & $0$ & $0$ & $0$ & $0$ & $0$\\
FreeSolv      & $0$ & $0$ & $0$ & $0$ & $0$ & $0$ & $0$ & $0$ & $0$\\
Lipophilicity & $0$ & $0$ & $0$ & $0$ & $0$ & $0$ & $0$ & $0$ & $0$\\
\bottomrule
\end{tabular}

\vspace{0.3em}
\noindent\small{Combined-diagnostic pass count of the Wuzi $100$-point grid across the full $\{0.90, 0.95, 0.99\} \times \{0.05, 0.10, 0.20\}$ threshold grid (nine $(\tau_r, \tau_{\mathrm{pcor}})$ combinations).}
\end{table}

\section{Feature-selection comparison and matched-model-class benchmark}

\begin{table}[h]
\centering
\small
\caption{Screen-pipeline configurations (RandomForest, rows from the RandomForest benchmark of the main text) vs.\ five external feature-selection comparators (LASSO, ElasticNet, mRMR, Boruta, $\ell_1$-Logistic) on the four MoleculeNet benchmarks. Regression metric is RMSE (lower is better); classification metric for BBBP is ROC-AUC (higher is better). Reported as mean across 5 folds with the same seed.}
\label{tab:fs_comparison}
\begin{tabular}{llrr}
\toprule
Dataset & Configuration & Metric & Mean $n$ features\\
\midrule
ESOL & \texttt{full} (RF)                 & RMSE $1.437$ & $30.0$\\
     & \texttt{pairwise\_pruned} (RF)     & RMSE $1.435$ & $\phantom{0}7.4$\\
     & \texttt{combined\_pruned} (RF)     & RMSE $1.475$ & $\phantom{0}4.4$\\
     & LASSO (linear, CV)                 & RMSE $1.545$ & $14.6$\\
     & ElasticNet (linear, CV)            & RMSE $1.546$ & $16.2$\\
     & mRMR $k{=}7$ (linear)              & RMSE $1.654$ & $\phantom{0}7.0$\\
     & Boruta + RF (linear scoring)       & RMSE $1.620$ & $10.0$\\
     & \emph{pairwise} $\to$ \emph{LASSO} (matched$^{\ddagger}$) & RMSE $1.580$ & $\phantom{0}7.2$\\
\midrule
FreeSolv & \texttt{full} (RF)             & RMSE $3.537$ & $30.0$\\
         & \texttt{pairwise\_pruned} (RF) & RMSE $3.540$ & $\phantom{0}8.0$\\
         & \texttt{combined\_pruned} (RF) & RMSE $3.561$ & $\phantom{0}3.2$\\
         & LASSO (linear, CV)             & RMSE $3.527$ & $13.4$\\
         & ElasticNet (linear, CV)        & RMSE $3.515$ & $18.6$\\
         & mRMR $k{=}8$ (linear)          & RMSE $3.618$ & $\phantom{0}8.0$\\
         & Boruta + RF (linear scoring)   & RMSE $3.588$ & $\phantom{0}2.0$\\
         & \emph{pairwise} $\to$ \emph{LASSO} (matched$^{\ddagger}$) & RMSE $3.465$ & $\phantom{0}6.4$\\
\midrule
Lipophilicity & \texttt{full} (RF)        & RMSE $1.030$ & $30.0$\\
              & \texttt{pairwise\_pruned} (RF) & RMSE $1.028$ & $\phantom{0}8.0$\\
              & \texttt{combined\_pruned} (RF) & RMSE $1.287^{\,\dagger}$ & $\phantom{0}0.2$\\
              & LASSO (linear, CV)        & RMSE $1.130$ & $15.6$\\
              & ElasticNet (linear, CV)   & RMSE $1.130$ & $17.8$\\
              & mRMR $k{=}8$ (linear)     & RMSE $1.142$ & $\phantom{0}8.0$\\
              & Boruta + RF (linear scoring) & RMSE $1.144$ & $\phantom{0}5.0$\\
              & \emph{pairwise} $\to$ \emph{LASSO} (matched$^{\ddagger}$) & RMSE $1.137$ & $\phantom{0}7.0$\\
\midrule
BBBP & \texttt{full} (RF)                 & AUC $0.846$ & $30.0$\\
     & \texttt{pairwise\_pruned} (RF)     & AUC $0.860$ & $\phantom{0}6.0$\\
     & \texttt{combined\_pruned} (RF)     & AUC $0.840$ & $\phantom{0}4.0$\\
     & $\ell_1$-Logistic (CV)             & AUC $0.790$ & $13.6$\\
     & mRMR $k{=}6$ (logistic)            & AUC $0.747$ & $\phantom{0}6.0$\\
     & Boruta + RF (logistic scoring)     & AUC $0.791$ & $23.0$\\
\bottomrule
\end{tabular}

\vspace{0.5em}
\noindent\small{$^\dagger$ Strict combined screen degenerates: retains zero features in $4/5$ folds; the RMSE shown is from the single fold that retained one feature. See the Limitations discussion of the main text.}
\noindent\small{$^\ddagger$ Pipeline-then-LASSO: in-fold pairwise screen $\to$ LASSO on the surviving features, same model class throughout. See this section and Table~\ref{tab:matched_class}.}
\end{table}

\begin{table}[!htbp]
\centering
\small
\caption{Pipeline-then-LASSO matched-model-class benchmark with paired $t$-tests. Same model class (LASSO for regression, $\ell_1$-Logistic for BBBP) and same 5-fold splits across all three configurations within each dataset, so the comparison isolates the effect of the upstream screening step. ``pre-LASSO feats'' is the mean number of features handed to LASSO after screening; ``final L1 feats'' is the mean number of features with non-zero LASSO coefficient. ``$t$ vs.\ full'' and ``$p$ vs.\ full'' are the paired $t$-statistic and two-sided $p$-value across the five folds (df $=4$) comparing the primary metric (RMSE for regression, ROC-AUC for BBBP) against the \texttt{full} (no-screen) row of the same dataset.}
\label{tab:matched_class}
\begin{tabular}{llcrrrr}
\toprule
Dataset & Configuration & Metric (mean $\pm$ std) & pre & final & $t$ vs.\ full & $p$ vs.\ full\\
\midrule
ESOL          & \texttt{full} (LASSO)                       & RMSE $1.536\pm 0.077$  & $30.0$         & $17.6$ & --- & ---\\
              & \texttt{pairwise} $\to$ LASSO               & RMSE $1.580\pm 0.085$  & $\phantom{0}7.4$ & $\phantom{0}7.2$ & $+3.77$ & $0.020^{*}$\\
              & \texttt{combined} $\to$ LASSO               & RMSE $1.588\pm 0.090$  & $\phantom{0}4.4$ & $\phantom{0}4.4$ & $+3.62$ & $0.022^{*}$\\
\midrule
FreeSolv      & \texttt{full} (LASSO)                       & RMSE $3.469\pm 0.400$  & $30.0$         & $12.8$ & --- & ---\\
              & \texttt{pairwise} $\to$ LASSO               & RMSE $3.465\pm 0.394$  & $\phantom{0}8.0$ & $\phantom{0}6.4$ & $-0.34$ & $0.748$\\
              & \texttt{combined} $\to$ LASSO               & RMSE $3.462\pm 0.387$  & $\phantom{0}3.2$ & $\phantom{0}3.2$ & $-0.46$ & $0.672$\\
\midrule
Lipophilicity & \texttt{full} (LASSO)                       & RMSE $1.128\pm 0.016$  & $30.0$         & $16.2$ & --- & ---\\
              & \texttt{pairwise} $\to$ LASSO               & RMSE $1.137\pm 0.013$  & $\phantom{0}8.0$ & $\phantom{0}7.0$ & $+3.45$ & $0.026^{*}$\\
              & \texttt{combined} $\to$ LASSO               & RMSE $1.220^{\,\dagger}$       & $\phantom{0}0.2$ & $\phantom{0}0.2$ & --- & ---\\
\midrule
BBBP          & \texttt{full} ($\ell_1$-Logistic)           & AUC $0.794\pm 0.019$   & $30.0$         & $11.6$ & --- & ---\\
              & \texttt{pairwise} $\to$ $\ell_1$-Logistic   & AUC $0.787\pm 0.026$   & $\phantom{0}6.0$ & $\phantom{0}6.0$ & $-1.71$ & $0.163$\\
              & \texttt{combined} $\to$ $\ell_1$-Logistic   & AUC $0.786\pm 0.028$   & $\phantom{0}4.0$ & $\phantom{0}4.0$ & $-1.56$ & $0.194$\\
\bottomrule
\end{tabular}

\vspace{0.5em}
\noindent\small{$^{*}$ significant at $\alpha=0.05$ (two-sided paired $t$-test, df $=4$). $^\dagger$ Combined screen retains zero features in $4/5$ Lipophilicity folds; the value shown is from the single fold that retained one feature. The same failure mode appears in the RandomForest benchmark of the main text and is treated as a formal limitation in the Limitations discussion of the main text.}
\end{table}

\section{Formal definitions of the 20 alternative-family candidate indices}\label{app:candidates}
% =====================================================================

For reproducibility we give the formal definition of each of the
twenty candidate indices screened in the candidate case study of the main text. Let
$G = (V, E)$ be a simple connected hydrogen-suppressed molecular
graph with $n = |V|$ vertices, $m = |E|$ edges, degree sequence
$\{d_v\}_{v\in V}$, geodesic distance matrix $D = (d_{uv})$,
adjacency matrix $A$, eccentricity $\varepsilon(v) = \max_u d_{uv}$,
and adjacency-spectrum / Laplacian-spectrum eigenvalues
$\{\lambda_i\}$ / $\{\mu_i\}$ respectively. All twenty definitions
are implemented in
\texttt{src/novel\_candidates.py} of the companion repository.

\paragraph{Information-theoretic (5).}
With $H(\{c_i\}) = -\sum_i (c_i / \sum_j c_j)\log(c_i / \sum_j c_j)$
the Shannon entropy of a count distribution:
\begin{align*}
\mathrm{InfoH\_deg}(G) &= H\bigl(\{|\{v : d_v = k\}|\}_k\bigr), \\
\mathrm{InfoH\_dist}(G) &= H\bigl(\{|\{(u,v) : d_{uv} = k,\, u<v\}|\}_k\bigr), \\
\mathrm{InfoH\_spec}(G) &= -\sum_{i} p_i \log p_i,\ p_i = \lambda_i^2 / \textstyle\sum_j \lambda_j^2, \\
\mathrm{InfoH\_ecc}(G) &= H\bigl(\{|\{v : \varepsilon(v) = k\}|\}_k\bigr), \\
\mathrm{Bonchev\_Id}(G) &= W\log_2 W - \sum_{k \ge 1} k\,d_k \log_2(k\,d_k),
\end{align*}
where in the last line $d_k = |\{(u,v) : d_{uv}=k,\,u<v\}|$ and
$W = \sum_k k\,d_k$ is the Wiener index.

\paragraph{Centrality-based (5).}
With standard NetworkX-style normalised centralities:
\begin{align*}
\mathrm{Btw\_sum}(G) &= \sum_{v \in V} c_B(v),\quad c_B(v) = \!\!\sum_{s \ne t \ne v}\!\!
\frac{\sigma_{st}(v)}{\sigma_{st}}, \\
\mathrm{Cls\_sum}(G) &= \sum_{v} (n-1)\Bigl(\sum_{u \ne v} d_{uv}\Bigr)^{-1}, \\
\mathrm{Eig\_sum}(G) &= \sum_{v} x_v,\ x \text{ the (positive) Perron eigenvector of } A, \\
\mathrm{Harm\_cent}(G) &= \sum_{v}\sum_{u \ne v}\frac{1}{d_{uv}}, \\
\mathrm{Btw\_var}(G) &= \mathrm{Var}_v\bigl(c_B(v)\bigr).
\end{align*}

\paragraph{Spectral-hybrid (3).}
With $\mathrm{diam}(G) = \max_{u,v} d_{uv}$ and Estrada index
$\mathrm{EE}(G) = \sum_i e^{\lambda_i}$:
\begin{align*}
\mathrm{AlgConn}(G) &= \mu_2(G), \text{ second-smallest Laplacian eigenvalue}, \\
\mathrm{Rho}\!\times\!\mathrm{Dia}(G) &= \rho(G)\cdot\mathrm{diam}(G),\ \rho(G) = \max_i|\lambda_i|, \\
\mathrm{EE}\!\times\!\mathrm{W}(G) &= \log\bigl(1+\mathrm{EE}(G)\bigr)\cdot\log\bigl(1+W(G)\bigr).
\end{align*}

\paragraph{Motif / structural (4).}
\begin{align*}
\mathrm{Triangles}(G) &= \tfrac{1}{6}\operatorname{tr}(A^3), \\
\mathrm{MeanClust}(G) &= \frac{1}{n}\sum_{v} \frac{2 t(v)}{d_v(d_v-1)},\ \,t(v) = \text{triangles at } v, \\
\mathrm{FourCyc}(G) &= \tfrac{1}{8}\bigl(\operatorname{tr}(A^4) - 2\!\sum_v d_v(d_v - 1) - 2m\bigr), \\
\mathrm{CutVerts}(G) &= |\{v \in V : G - v \text{ is disconnected}\}|.
\end{align*}

\paragraph{Eccentricity-based (3).}
\begin{align*}
\mathrm{TotEcc}(G) &= \sum_{v \in V} \varepsilon(v), \\
\mathrm{Radius}(G) &= \min_{v \in V} \varepsilon(v), \\
\mathrm{MeanEcc}(G) &= \frac{1}{n}\sum_{v \in V} \varepsilon(v).
\end{align*}

None of these twenty candidates is claimed as a newly invented
index. The information-theoretic family follows Bonchev and
Trinajsti\'c~\cite{Gutman2013Degree, TodeschiniConsonniHandbook};
the centrality family follows the standard graph-theory
literature; the motif and eccentricity families are textbook
graph invariants; the spectral-hybrid combinations are
deliberately constructed novel pairings to test whether
non-standard product indices escape the redundancy ceiling. They
appear here purely as screening targets to test the generality
of the combined diagnostic of the screening protocol of the main text, and the
verdict reported in the candidate case study of the main text is a statement about
this particular candidate set under the tested baselines and
thresholds, not a claim that any individual candidate is or is
not a useful QSPR descriptor.

% =====================================================================

\begin{thebibliography}{9}\small
\bibitem{Gutman2013Degree} I.~Gutman, \emph{Degree-based topological indices}, Croat.~Chem.~Acta \textbf{86}(4) (2013), 351--361.
\bibitem{TodeschiniConsonniHandbook} R.~Todeschini and V.~Consonni, \emph{Molecular Descriptors for Chemoinformatics}, 2nd ed., Wiley-VCH, Weinheim, 2009.
\end{thebibliography}
\end{document}
